\section{Organizations}
\label{sec:organizations}

This section asks what the platform (\cref{sec:platform}) could offer an organization: World B,
applied to a firm's own artifacts instead of public text. The PoC evidence behind this section is
the N3 gap map (\cref{sec:n3}), run on two demo domains built entirely from \emph{public} text: 35
fetched PLC-troubleshooting web pages (forums and vendor blogs, with 42 promotional seeds excluded),
and the Article~29 Working Party's WP248 guidance on data protection impact assessments (GDPR v1,
whose support is 80\% guidance from the UK Information Commissioner's Office (ICO) and statute, across three independence keys). Its own
closeout reports that the closure step is not informative against a mismatched-evidence control on either domain,
and that one adversarial AI reviewer judged 3 of 19 candidates precise (strict reading) on the pre-final maps, in-sample and with no human coder. Nothing below claims the PoC has run on real organizational
data, and nothing below claims it identifies real tacit knowledge. \evint\ for every application
below, unless marked otherwise.

\subsection{Three organizational blind spots}
\label{sec:org-blindspots}

\textbf{A technician's diagnostic cue.} Orr's ethnography of photocopier repair technicians found
that diagnosis proceeds through a narrative process: technicians build a coherent story of the
failing machine, and that story, not the service manual, keeps the group's diagnostic knowledge
current \parencite{orr1996talking}. Brown and Duguid read the same fieldwork as knowledge held by a
community of practice, not by individuals or documents \parencite{brown1991organizational}.
\evlit. The PLC demo produced a candidate of the same shape: a record built on the span ``a common
mistake\dots{} is panicking and abandoning systematic approaches,'' where the absence check found
at most a partial statement of how a troubleshooter notices the drift (state: \code{partial}).
\evobs\ (candidate, not validated) that a lens can flag this \emph{kind} of gap in text; whether it
is real tacit knowledge is untested.

\textbf{A senior engineer's design rationale.} Szulanski's study of 122 practice transfers across
eight firms found that the major barriers to transfer were knowledge-related: the recipient's lack of
absorptive capacity, \emph{causal ambiguity} (not knowing why a practice works) and an arduous
relationship between source and recipient, rather than mainly motivational factors \parencite{szulanski1996exploring}. \evlit. In software,
this shows up as commits and configuration changes with no issue, pull request, or design record
explaining the reasoning. \evint.

\textbf{A compliance officer's exception judgment.} Hollnagel's work-as-imagined versus
work-as-done distinction \parencite{hollnagel2018safety} is the general frame: the written rule is
the imagined process, and a practitioner's case-by-case judgment of borderline cases is the work
actually done. \evlit. WP248 is written in exactly this shape, requiring a data protection impact
assessment ``in most cases'' when two screening criteria are met and ``in some cases'' when only
one is, without saying which. The GDPR demo's best-rated candidate is built on precisely this
hedge. \evobs\ (candidate, not validated) --- the run is not N3-admissible (140 pending verdicts),
and the closure check that would confirm the boundary is genuinely unstated elsewhere is, again,
uninformative against its mismatched-evidence control.

\subsection{World B evidence versus World A}
\label{sec:org-worldb}

World A evidence (public text) is expert-voiced explanation, written to be read. World B evidence
is organizational trace: SOPs and runbooks (\code{imagined}), but also tickets, incident reports,
commit history, review threads, chat, and postmortems (\code{done}). The doc-versus-practice
divergence is a World-B-specific signal public text cannot show at all, having no paired execution
log to compare against \parencite{hollnagel2018safety}. \evlit\ for the concept; \evfut\ for a
benchmark that would prove a detector can find it reliably.

World B differs from World A in three ways the platform's design already treats separately
(\cref{sec:platform-nfr}):
\begin{itemize}[noitemsep]
  \item \emph{Access control}: retrieval must enforce per-user rights, and the standard
    retrieval-augmented-generation evaluation frameworks we found
    \parencite{gao2023retrieval,es2024ragas} do not score access-scoped retrieval against a
    held-out knowledge criterion \evhyp.
  \item \emph{Confidentiality and intellectual property}: sending artifacts to a third-party model
    is the organization's decision, not the pipeline's.
  \item \emph{Personal-data exposure}: Git authorship and review participation are the raw material
    for scoring who alone understands a component, which is personal-data processing about
    identifiable employees (\cref{sec:org-ethics}).
\end{itemize}

\subsection{How the pipeline changes for a private corpus}
\label{sec:org-pipeline}

The reconstruction-to-gap-map pipeline's (\cref{sec:approach}) core components are domain-neutral
by construction: they live in the core evidence package, independent of any one domain (decision
0006). The lexicons, denylist, and stop-anchors that drive the lenses are not: they are English and
were tuned in-sample on PLC and GDPR (\cref{sec:n3}), so a World B run needs the same in-sample
re-tuning risk any new domain carries, not a domain-neutral starting point. Applying the pipeline to
World B needs three further changes, none built or tested. \evfut{} throughout.

\textbf{Retrieval over a private, access-scoped corpus.} Reconstruction currently retrieves the open
web; a World B run needs an indexed corpus (SOPs, tickets, Git history, design records) with
per-document access rights carried through to every claim's provenance. \evhyp.

\textbf{On-premises or contractually isolated models.} Confidentiality likely forecloses sending
artifacts to a third-party API by default. The gap map's own closure judge already runs locally,
though the recorded reason was model-family separation and zero cost, not confidentiality (decision
0008 (proposed); \cref{sec:n3}) --- a precedent that on-premises deployment is viable, not that it
works on organizational text: its closure judgments were uninformative against a mismatched-evidence control on
either PoC demo domain. \evobs\ for viability; \evint\ for organizational applicability.

\textbf{Gap map, question ranking, and residual capture, unchanged in shape.} Lenses fire on
constructs the corpus attests but does not state; a judge assesses whether the gap is genuinely
unclosed elsewhere; the ranked output becomes questions for a named or pseudonymized expert.
Untested for World B specifically: whether lenses built for expert prose fire usefully on ticket
and commit text, and whether closure against a private corpus is any more informative than the
public-text step that failed its control on both PoC domains.

\subsection{Value propositions, stated soberly}
\label{sec:org-value}

\Cref{tab:org-value} separates what the literature establishes about the underlying problem from
what is actually shown about this system solving it. No credible, peer-reviewed cost figure for
organizational knowledge loss in general was found in the evidence review behind this report, and
none is asserted here.

\begin{table}[htbp]
  \centering
  \footnotesize
  \caption{Organizational value propositions: literature grounding versus what this system has
    shown.}
  \label{tab:org-value}
  \begin{tabularx}{\linewidth}{@{}p{0.19\linewidth}Xp{0.30\linewidth}@{}}
    \toprule
    Proposition & Literature grounding & Evidence that our system solves it \\
    \midrule
    Onboarding acceleration &
      A taxonomy of documentation issues built from 878 software artifacts
      \parencite{aghajani2019software}; documentation problems are one of five barrier categories
      for newcomers to open source projects \parencite{steinmacher2014barriers}; we found no study
      of onboarding speed &
      \evhyp\ Untested that a gap map shortens onboarding \\
    Retiring-expert knowledge capture &
      Retention frameworks treat this as a real, prioritized risk, on case-study and practitioner
      argument \parencite{delong2004lost} &
      \evhyp\ No validated method converts an artifact gap map into a capture priority \\
    Turnover-induced knowledge loss, quantified &
      Established at scale in software: 65\% of 133 popular GitHub projects have a truck factor of
      2 or fewer \parencite{avelino2016novel}; in Chrome and a project at Avaya, the mean of the
      worst 5\% of quarterly knowledge losses (abandoned files) was 3.6--3.8 times the expected
      loss, and historical simulations of departures gave losses over five times the expected loss
      \parencite{rigby2016quantifying}, and a
      replication found more severe extremes \parencite{nassif2017revisiting}; 315 of 1{,}932
      projects (16\%) were abandoned after truck-factor developers left
      \parencite{avelino2019abandonment} &
      \evlit\ that the risk is real and measurable in software. \evhyp\ that an artifact-coverage
      score adds predictive value beyond authorship-only concentration --- exactly what the
      program's H-OSS study (V3) is designed to test (\cref{subsec:hoss}), and it has not run \\
    Troubleshooting and diagnostic support &
      Diagnostic knowledge circulates as narrative among practitioners, not in manuals
      \parencite{orr1996talking} &
      \evhyp\ Untested whether reconstruction plus a gap map surfaces the missing cue rather than
      just the documented procedure \\
    Audit and regulatory defensibility &
      The gap between work as imagined and work as done has a theoretical frame
      \parencite{hollnagel2018safety}, and conformance checking compares event logs with process
      models \parencite{carmona2018conformance}; our search found no public corpus joining SOPs to
      execution logs &
      \evfut\ We found no benchmark to validate an SOP-versus-record divergence detector against \\
    \bottomrule
  \end{tabularx}
\end{table}

Open-source software is also the one organizational domain where knowledge loss can be measured
without any interview at all. The domain strategy (D2, \cref{sec:domains}) therefore uses it for the
program's zero-participant organizational test, H-OSS (\cref{subsec:hoss}). H-OSS is a separate
hypothesis: its result can neither support nor refute H2.

\subsection{A candidate organizational pilot}
\label{sec:org-pilot}

\textbf{Partner type (Recommended).} A moderate-size software organization with an internal Git
repository, issue tracker, code review history, and some design records, closest to the program's
registered H-OSS design (V3, \cref{subsec:hoss}), reusing a pre-specified analysis rather than
inventing one. An industrial-maintenance partner is the stronger long-run target for the
diagnostic-cue value proposition, but lacks as clean a natural experiment as developer departure and
tends toward terse, jargon-heavy text, a documented low-support condition; start with the software
partner and treat maintenance as a second-wave target.

\textbf{Domain, data, and measures.} Git history, code review comments, the issue tracker, and
design records, all dated; no private chat without separate, explicit consent. In the program this
pilot is V7: it starts only after MS5 reads continue and the deferred N2 validation is complete
(\cref{sec:validation}). The procedure has the same structure as H-OSS:
\begin{enumerate}[noitemsep]
  \item freeze the corpus shortly before a known or planned senior-contributor departure;
  \item compute the gap map and a concentration score per component from artifacts dated at or
    before that freeze;
  \item after the departure, blind-code post-freeze rationale-seeking issues and defect-fix
    commits per component;
  \item compare the ranking against the strongest baseline (truck factor, code churn,
    documentation link density).
\end{enumerate}

\textbf{Success and failure criteria.} Success: the ranking predicts post-departure trouble better
than the strongest baseline, by a pre-registered margin, mirroring the H-OSS gate. Failure: it
does not, meaning artifact coverage adds nothing beyond authorship concentration --- informative,
not wasted, mirroring N3's own honest closure-check failure (\cref{sec:n3}).

\subsection{Ethics and data protection}
\label{sec:org-ethics}

The points in this subsection follow the reorientation plan's own reading of applicable law and
guidance. \textbf{They are stated as information gathered for planning purposes, not as
independently verified legal advice, and a qualified data protection counsel should review them
before any pilot begins.}

A data protection impact assessment is required wherever per-person concentration scoring performs
systematic processing of identifiable employees' work. The proposed legal basis is legitimate
interest with a documented balancing test, not consent, since consent is generally invalid inside an
employment relationship given the power imbalance it carries. Where a German entity is involved, a
works council agreement is required before deployment: a system objectively capable of assessing
employee performance is subject to co-determination under German works constitution law regardless
of the deployer's intent. The pilot must not feed task allocation, promotion, or performance
evaluation, staying clear of the stricter category high-risk employment AI systems fall into under
EU law. Outputs default to aggregate or pseudonymized scores, named only by consent. Company
artifacts remain the firm's intellectual property; a written agreement should bar training any model
outside the engagement and require on-premises processing or an equivalent data-processing
agreement.

\subsection{Risks specific to organizations}
\label{sec:org-risks}

\textbf{Goodhart's law applied to documentation.} If a coverage score becomes a target metric, the
response most available to people is to write documentation that raises the score, not
documentation that conveys the missing knowledge. A large empirical taxonomy of software documentation
issues, built from 878 artifacts in mailing lists, Stack Overflow, issue trackers and pull requests,
shows that documentation problems arise without any scoring pressure
\parencite{aghajani2019software}. \evlit\ for the taxonomy; \evint\ that a coverage
key-performance-indicator would plausibly make by-the-letter-but-hollow documentation more common.

\textbf{Surveillance framing.} Szulanski's evidence runs against the intuitive story that experts
withhold knowledge as leverage: the major barriers it found were knowledge-related (absorptive
capacity, causal ambiguity, an arduous source--recipient relationship) rather than mainly motivational
\parencite{szulanski1996exploring}. \evlit. The more defensible risk is that a
system inferring ``only one person understands this'' from Git and review history is doing exactly
the processing the legitimate-interest test and works-council co-determination exist to check.
\evint. Deployed without \cref{sec:org-ethics}'s safeguards, it reads to staff as covert
performance monitoring, which alone can block adoption or trigger a works-council veto.

\textbf{Confidentiality and injection exposure.} Every internal document a retrieval system serves
back is a channel to plant instructions; indirect prompt injection through retrieved content is a
demonstrated attack class \parencite{greshake2023not}. \evlit. A World B deployment retrieving
tickets, chat, and commit messages has a larger, less curated attack surface than a World A pipeline
reading public papers. \evint. The platform's answer is that retrieved text is always data, never
instructions (\cref{sec:platform-nfr}).

\textbf{Misuse of concentration scores.} A per-person ``irreplaceability'' score is a short step
from a redundancy-targeting or performance tool, the regulatory trigger \cref{sec:org-ethics} is
built to avoid; the mitigation (aggregate by default, consent for naming, a contractual ban on
evaluation use) is a governance decision enforced by contract, since the pipeline cannot enforce it
itself. \evint.
